% Necoyoad Architecture Blueprint v10 — Technical Assessment and Remastering Feasibility
% This volume is PRESCRIPTIVE (unlike v1-v9 which were descriptive)

\documentclass[11pt,a4paper]{article}
\usepackage[T1]{fontenc}
\usepackage[utf8]{inputenc}
\usepackage{lmodern}
\usepackage{microtype}
\usepackage{graphicx}
\usepackage{xcolor}
\usepackage{geometry}
\usepackage{amsmath}
\usepackage{amssymb}
\usepackage{tcolorbox}
\tcbuselibrary{skins,breakable,listings}
\usepackage{colortbl}
\usepackage{booktabs}
\usepackage{tabularx}
\usepackage{longtable}
\usepackage{array}
\usepackage{enumitem}
\usepackage{adjustbox}
\usepackage{caption}
\usepackage{fancyhdr}
\usepackage{titlesec}
\usepackage{pdflscape}
\usepackage{listings}
\usepackage{tikz}
\usetikzlibrary{arrows.meta,positioning,calc,shapes.geometric,fit,backgrounds,chains,decorations.pathreplacing,shapes.multipart}
\usepackage[colorlinks=true,linkcolor=blue!50!black,citecolor=blue!40!black,urlcolor=blue!60!black,bookmarks=true,bookmarksnumbered=true,unicode=true,pdftitle={Necoyoad - Technical Assessment and Remastering Feasibility (v10)},pdfauthor={Reverse-Engineered Architecture Report},pdfsubject={Technical Assessment, Code Review, Remastering Feasibility, Laravel Migration},pdfkeywords={Necoyoad, OpenCart, PHP 8.1, Technical Assessment, Code Review, Remastering, Laravel, Symfony, Migration, Feasibility}]{hyperref}
\geometry{a4paper, top=2.4cm, bottom=2.4cm, left=2.4cm, right=2.4cm}
\setlength{\headheight}{14pt}
\pagestyle{fancy}
\fancyhf{}
\fancyhead[L]{\small\itshape Necoyoad v10 --- Technical Assessment \& Remastering Feasibility}
\fancyfoot[C]{\small\thepage}
\renewcommand{\headrulewidth}{0.3pt}
\renewcommand{\footrulewidth}{0pt}
\titleformat{\section}{\Large\bfseries\color{black!85}}{\thesection}{0.6em}{}
\titleformat{\subsection}{\large\bfseries\color{black!80}}{\thesubsection}{0.6em}{}
\titleformat{\subsubsection}{\normalsize\bfseries\color{black!75}}{\thesubsubsection}{0.5em}{}
\widowpenalty=10000
\clubpenalty=10000
\emergencystretch=4em
\hbadness=10000
\hfuzz=2pt
\vfuzz=2pt
\sloppy

\definecolor{necoyoad-blue}{HTML}{1F4E79}
\definecolor{necoyoad-orange}{HTML}{B85C00}
\definecolor{necoyoad-green}{HTML}{2E6B33}
\definecolor{necoyoad-red}{HTML}{B22222}
\definecolor{nodefill}{HTML}{F4F7FB}
\definecolor{nodefillalt}{HTML}{FBF4EE}
\definecolor{nodefillgreen}{HTML}{F2F8F2}
\definecolor{nodefillred}{HTML}{FFF0F0}
\definecolor{nodefillgray}{HTML}{F0F0F2}

\newtcolorbox{goodbox}[1][]{
  enhanced, breakable,
  colback=nodefillgreen, colframe=necoyoad-green!70,
  boxrule=0.4pt, arc=2pt, left=8pt, right=8pt, top=6pt, bottom=6pt,
  fonttitle=\bfseries\small, coltitle=white,
  title=#1
}
\newtcolorbox{badbox}[1][]{
  enhanced, breakable,
  colback=nodefillred, colframe=necoyoad-red!70,
  boxrule=0.4pt, arc=2pt, left=8pt, right=8pt, top=6pt, bottom=6pt,
  fonttitle=\bfseries\small, coltitle=white,
  title=#1
}
\newtcolorbox{reusebox}[1][]{
  enhanced, breakable,
  colback=nodefill, colframe=necoyoad-blue!60,
  boxrule=0.4pt, arc=2pt, left=8pt, right=8pt, top=6pt, bottom=6pt,
  fonttitle=\bfseries\small, coltitle=white,
  title=#1
}

\captionsetup{justification=raggedright, singlelinecheck=false, font=small, labelfont=bf}

\begin{document}
\thispagestyle{empty}
\begin{center}
{\Large\bfseries Necoyoad --- Technical Assessment \& Remastering Feasibility (v10)}\\[0.4em]
{\small\itshape What was good, what was bad, what to keep, and whether to rebuild\\
from scratch with Laravel or another modern framework}\\[1.2em]
\end{center}

\begin{abstract}
\noindent
This is the tenth and final volume of the Necoyoad architecture blueprint.
Unlike volumes v1--v9 (which were strictly descriptive), v10 is
\textbf{prescriptive}: it evaluates the codebase, identifies what was well
designed, what was poorly designed, what is worth preserving in a remastering
effort, and assesses the feasibility of rebuilding the platform from scratch
using Laravel (or a similar modern PHP framework).

The assessment is based on the exhaustive source-code analysis conducted
across v1--v9: 1,450 PHP files, 87 database tables, 202 pages of verified
architecture documentation. The findings are organised into four parts:
\textbf{Part I --- What Was Good} (10 architectural strengths worth
preserving); \textbf{Part II --- What Was Bad} (15 technical debts and
design flaws); \textbf{Part III --- What to Keep} (8 patterns and
abstractions that should survive a remastering); and \textbf{Part IV ---
Feasibility of a Laravel Rebuild} (a pattern-by-pattern mapping from
Necoyoad to Laravel, with effort estimates and a recommended migration
strategy).
\end{abstract}

\vspace{0.6em}
\noindent\textbf{Keywords:} Necoyoad, OpenCart, PHP 8.1, technical
assessment, code review, remastering, Laravel, Symfony, migration
feasibility, architecture patterns.

\newpage
\tableofcontents
\newpage

% ============================================================
\section{Introduction}\label{sec:intro}
% ============================================================

\subsection{From Descriptive to Prescriptive}

Volumes v1--v9 were strictly descriptive: they documented how the system
\emph{is}, without proposing improvements. This volume changes stance.
The user has asked for an evaluative assessment: what was good, what was
bad, what can be reused, and whether to rebuild from scratch. v10 answers
those questions directly.

\subsection{Basis for Assessment}

The assessment draws on:
\begin{itemize}[leftmargin=1.6em,itemsep=2pt]
  \item v1 (SQL-only reconstruction) --- the 87-table schema.
  \item v2 (source verification) --- the bootstrap, engine, auth, API.
  \item v3 (rendering pipeline) --- the widget composition system.
  \item v4 (marketing subsystem) --- the campaign pipeline.
  \item v5 (multi-store/multi-language) --- the polymorphic object spine.
  \item v6 (admin back-office) --- the AdminController declarative CRUD.
  \item v7 (menu links) --- the recursive tree rendering.
  \item v8 (CMS posts/pages) --- the widget composition modes.
  \item v9 (banner subsystem) --- the jQuery plugin discriminator.
\end{itemize}

Total: 216 pages of verified architecture documentation, covering every
subsystem of the platform.

\subsection{Context}

Necoyoad is a PHP 8.1 / MySQL 5.7 e-commerce and CMS platform, derived
from OpenCart 1.x/2.x, with significant proprietary extensions. The
codebase is $\sim$17,000 lines of PHP (excluding vendor libraries),
1,450 files, 87 database tables. It was built for the Venezuelan market
(RIF tax IDs, bank-transfer payments, Spanish-first UI) circa 2009--2023.

% ============================================================
\part{What Was Good}\label{part:good}
% ============================================================

\section{Architectural Strengths}\label{sec:good}

\subsection{The AdminController Declarative CRUD Framework}

\begin{goodbox}[AdminController --- the crown jewel of the admin]
The \texttt{AdminController} base class (1,244 lines, v6) is a declarative
CRUD framework that reduces 70+ admin controllers to a few lines of
configuration each. A subclass declares:
\texttt{\$form\_vars} (field definitions with types),
\texttt{\$filters} (list-view filter definitions),
\texttt{\$public\_methods} (accessible CRUD methods),
\texttt{\$model\_name}/\texttt{\$model\_route} (the model to load),
and \texttt{\$object\_type} (the polymorphic discriminator).

The base provides the entire CRUD lifecycle:
\texttt{index()} (list), \texttt{insert()}/\texttt{update()} (via
shared \texttt{upsert()}), \texttt{copy()}, \texttt{delete()},
\texttt{activate()}, \texttt{sortable()}, \texttt{grid()} (AJAX grid),
\texttt{getList()}, \texttt{getForm()}, \texttt{validateForm()},
\texttt{validateDelete()}.

This is a significant abstraction over vanilla OpenCart (where each
controller has its own CRUD) and over many modern frameworks (where
CRUD is typically generated by Artisan commands rather than declared
at the class level). The declarative approach means adding a new admin
entity is a matter of writing a controller with a few properties, not
copy-pasting 500 lines of boilerplate.
\end{goodbox}

\subsection{The Polymorphic Object Spine}

\begin{goodbox}[Five-table polymorphic spine --- scalable and extensible]
The five polymorphic tables (\texttt{object\_to\_store},
\texttt{object\_to\_category}, \texttt{description},
\texttt{property}, \texttt{url\_alias}) use the
\texttt{(object\_id, object\_type)} pair as a composite soft foreign
key. This means:

\begin{itemize}[leftmargin=1.6em,itemsep=2pt]
  \item Adding a new content type (e.g.\ ``survey'', ``auction'')
        requires \emph{no schema changes} --- just populate
        \texttt{object\_type} with the new string.
  \item Localisation is built-in: every content type gets per-language
        descriptions via the \texttt{description} table's
        \texttt{(object\_id, object\_type, language\_id)} triple.
  \item SEO URLs are built-in: every content type gets per-language
        slugs via \texttt{url\_alias}.
  \item Free-form metadata is built-in: every content type gets EAV
        properties via the \texttt{property} table.
  \item Multi-store scoping is built-in: every content type gets
        store assignment via \texttt{object\_to\_store}.
\end{itemize}

This is a well-designed polymorphic architecture that anticipates
future content types. It is more flexible than Laravel's polymorphic
relations (which require a morph map) because the discriminator is a
free-form string, not a class name.
\end{goodbox}

\subsection{The Dual Hooks/Events Extension System}

\begin{goodbox}[Hooks + Events --- the right tool for each job]
The dual extension system (v2, v3) provides:

\begin{itemize}[leftmargin=1.6em,itemsep=2pt]
  \item \textbf{Hooks} (WordPress-style): \texttt{addFilter},
        \texttt{applyFilters}, \texttt{addAction}, \texttt{run}.
        Filters chain values through callbacks; actions short-circuit
        on truthy return. Used for \texttt{render}, \texttt{fetch},
        \texttt{loadWidgets}, \texttt{csspath}, \texttt{module:settings}.
  \item \textbf{Events} (Node-style): \texttt{on}, \texttt{emit},
        \texttt{off}. No short-circuit; all listeners fire. Used for
        \texttt{dispatch}, \texttt{beforeLoad}, \texttt{afterLoad},
        \texttt{registry:update}, \texttt{php:error}.
\end{itemize}

The separation is architecturally sound: Hooks are for
\emph{transforming} values and \emph{short-circuiting} flows (I want
to change the rendered HTML); Events are for \emph{broadcasting}
occurrences (I want to know when a dispatch happens). Many modern
frameworks conflate these into a single event system, which leads to
confusion about whether a listener can modify the flow.
\end{goodbox}

\subsection{The Widget Composition System}

\begin{goodbox}[Three-mode widget composition --- flexible and powerful]
The widget composition system (v3, v8) supports three modes:

\begin{enumerate}[leftmargin=1.6em,itemsep=2pt]
  \item \textbf{Dynamic}: admin configures widgets via the visual
        editor; \texttt{widgets-rows.tpl} emits
        \texttt{\{\%widget\_name\%\}} tokens; \texttt{fetch()}
        substitutes them.
  \item \textbf{Manual}: template author hardcodes
        \texttt{\{\%widget\_name\%\}} tokens in a \texttt{.tpl} file.
  \item \textbf{Hybrid}: combines dynamic sections with hardcoded
        sections.
\end{enumerate}

Plus: per-entity widget overrides (product-specific widgets),
per-item widget composition (banner slide overlays), the
\texttt{only:} prefix for embedded pages, and the
\texttt{widgets-common.tpl} universal layout. This is a page-builder
experience that rivals WordPress's block editor or Drupal's Layout
Builder, implemented in raw PHP.
\end{goodbox}

\subsection{The Per-Entity Template Override}

\begin{goodbox}[Per-entity template override via EAV]
Each individual post, page, product, or category can use a different
template, set via the admin's visual editor and stored in the EAV
\texttt{property('style', 'view')} key. The resolution is 3-level:
EAV property $\to$ config default $\to$ hardcoded fallback. This
allows a blog post to use a standard template, a landing page to use
a custom template with SVG sections, and a product page to use a
custom product-info template --- all without code changes.
\end{goodbox}

\subsection{The Banner jQuery Plugin Discriminator}

\begin{goodbox}[jquery-plugin discriminator --- plugin-aware widget]
The \texttt{banner.jquery\_plugin} column (v9) drives the entire
rendering pipeline: template selection, JS enqueuing, CSS enqueuing,
and the \texttt{ntPlugins} registry push. This is the only widget
module that dynamically selects its template from a database column.
The 33-template library (NivoSlider, Slick, Camera, Fancybox, etc.)
with per-template JS/CSS assets is a rich visual content system.
\end{goodbox}

\subsection{The Multi-Store / Multi-Language Composition}

\begin{goodbox}[Store x Language matrix --- data-level multi-tenancy]
The multi-store architecture (v5) is a clean data-level multi-tenancy:
\texttt{STORE\_ID} as a global constant, filtered at query time by
every model. The 6-level language detection priority chain
(GET $\to$ session $\to$ cookie $\to$ browser $\to$ config) is thorough.
The $S \times L$ composition matrix (5 filtering dimensions at query
time) produces a unique content surface per (store, language) pair
without a dedicated schema construct.
\end{goodbox}

\subsection{The Campaign Marketing Pipeline}

\begin{goodbox}[Self-hosted Mailchimp clone]
The marketing subsystem (v4) is a complete email-marketing pipeline:
newsletter composition with WYSIWYG, contact-list targeting,
scheduled sending via cron, per-recipient personalisation (7 tokens +
13 store-level tokens), open/click tracking with full HTTP context,
and a 50-email-per-pass throttle with 15-minute defer. This is more
than many commercial e-commerce platforms offered in 2009, and it is
fully self-hosted.
\end{goodbox}

\subsection{The Visual Theme Editor}

\begin{goodbox}[Visual theme editor with nt-editable hooks]
The visual theme editor (v3) is built into the HTML itself via
\texttt{nt-editable}, \texttt{data-widget}, \texttt{data-position},
\texttt{data-necotienda\_module} attributes. The editor's JS scans
the DOM and attaches drag/click handlers. The
\texttt{/**key**/ ... /**/key**/} CSS markers allow per-widget CSS
editing. The \texttt{theme\_style} table stores CSS-rule overrides
at the selector/property/value granularity. This is a
no-code customisation experience.
\end{goodbox}

\subsection{The \texttt{ntPlugins} JavaScript Registry}

\begin{goodbox}[ntPlugins --- decoupled deferred initialisation]
The \texttt{ntPlugins} registry (v9) is a decoupled, deferred-
initialisation pattern for jQuery plugins. Each slider template
pushes a \texttt{\{id, config, plugin\}} entry, and
\texttt{loadNTPlugins()} applies the plugin on DOM ready. This allows
multiple sliders on the same page without templates knowing about
each other --- a clean observer pattern for the client side.
\end{goodbox}

% ============================================================
\part{What Was Bad}\label{part:bad}
% ============================================================

\section{Technical Debt and Design Flaws}\label{sec:bad}

\subsection{MyISAM Engine (No Transactions, No FKs)}

\begin{badbox}[MyISAM --- the root cause of many issues]
All 87 tables use MyISAM. This means: no transactions (multi-table
writes like order creation are not atomic), no foreign keys (all
\texttt{\_id} columns are soft references, referential integrity is
the application's responsibility), table-level locking (concurrent
writes to \texttt{stat}, \texttt{campaign\_stat}, \texttt{search}
block each other), and no crash recovery. This is the single
biggest technical debt in the codebase.
\end{badbox}

\subsection{Double-MD5 Password Hashing}

\begin{badbox}[Double-MD5 with per-row salt --- insecure by modern standards]
The password algorithm (v2) is double-MD5 with a per-row random salt:
\texttt{md5(md5(plaintext) . suffix) . ':' . suffix}. While the
salt prevents rainbow-table attacks, MD5 is computationally cheap
(making brute-force feasible) and the double-MD5 offers no meaningful
security improvement over single-MD5. Modern standards require
bcrypt, argon2id, or scrypt --- slow hashes designed to make
brute-force infeasible.
\end{badbox}

\subsection{The \texttt{LIKE '\%key=value\%'} Widget Query Pattern}

\begin{badbox}[LIKE pattern --- full-table scans on every widget lookup]
Every widget filter (v3, v5) is a \texttt{LIKE} query against the
serialised settings blob: \texttt{WHERE p.value LIKE ...}. This forces a full-table scan on every
widget lookup. On a page with 5 layout positions, each with 3 rows,
each with 2 columns, each with 4 widgets, that is 165 full-table
scans. The cache absorbs most of the cost, but the first request (or
any request after a cache clear) is extremely slow.
\end{badbox}

\subsection{N+1 Query Patterns}

\begin{badbox}[Recursive query explosion in widget and menu trees]
The recursive \texttt{getLinks()} method (v7) and the per-object
widget query (v5) both exhibit N+1 patterns: one query per parent,
one query per child, one query per widget. A 3-level menu with 5/3/2
items generates 63 queries on cache miss. The widget system's
two-query merge (default tree + per-object tree) doubles the query
count for every layout position.
\end{badbox}

\subsection{No API Authentication}

\begin{badbox}[validateTokens() stub --- the API is fully open]
The admin API's \texttt{validateTokens()} method (v2) is a stub that
always returns \texttt{true}. The TODO comment acknowledges that
public/private API keys, tokens, expiry, access level, licensing, and
premium membership checks were planned but never implemented. The
only protection is the admin session token pre-action. Any caller
with an admin session can hit every API endpoint.
\end{badbox}

\subsection{Vestigial Bounce Processing}

\begin{badbox}[Bounce processor is dead code]
The bounce processor (v4) is commented out of the cron and references
non-existent Interspire SendStudio classes. There is no working
bounce detection. Bounced emails are not automatically unsubscribed;
the contact stays active and will be emailed again on the next
campaign.
\end{badbox}

\subsection{Inconsistent \texttt{object\_to\_store} Writes}

\begin{badbox}[Legacy and polymorphic columns coexist in the same table]
The admin store model (v5) writes products, categories, and
manufacturers using legacy column names (\texttt{product\_id},
\texttt{category\_id}, \texttt{manufacturer\_id}), but writes
downloads, pages, and banners using the polymorphic pair
(\texttt{object\_id}, \texttt{object\_type}). The storefront reads
using the polymorphic pair. Posts and post-categories \texttt{INSERT}
statements omit \texttt{store\_id} entirely. This inconsistency means
some content types are not correctly store-scoped.
\end{badbox}

\subsection{Broken API PUT Branches}

\begin{badbox}[All 5 marketing API PUT branches are broken]
The PUT (update) branches of all 5 marketing API endpoints (v4) use
the loaded DB row as the data source instead of the JSON body, or
reference an undefined \texttt{\$query} variable (mailservers). The
GET, POST, and DELETE branches work. Only PUT is broken.
\end{badbox}

\subsection{PHPMailer 5.0 from 2009}

\begin{badbox}[PHPMailer 5.0 --- 14 years old and insecure]
The campaign sender (v4) uses PHPMailer 5.0 (released 2009), renamed
\texttt{Mailer}. It supports only \texttt{AUTH LOGIN} (no OAuth2, no
CRAM-MD5). Multiple CVEs have been disclosed against PHPMailer 5.x
over the years. The SMTP client is a from-scratch implementation,
not a wrapper around PHP's \texttt{mail()}.
\end{badbox}

\subsection{No Unsubscribe Mechanism}

\begin{badbox}[No unsubscribe --- CAN-SPAM/GDPR non-compliance]
There is no \texttt{?unsubscribe=1} endpoint and no
\texttt{List-Unsubscribe} email header (v4). Recipients cannot
opt out of campaigns. This is non-compliant with CAN-SPAM (US),
GDPR (EU), and most modern email regulations.
\end{badbox}

\subsection{Route-Level RBAC Only}

\begin{badbox}[No object-level permissions]
The admin RBAC (v6) checks whether the route (e.g.\ \texttt{store/product})
is in the user group's \texttt{access}/\texttt{modify} arrays. It does
not check whether the user has permission to edit a \emph{specific}
product. Any admin with \texttt{modify} permission on
\texttt{store/product} can edit any product in any store.
\end{badbox}

\subsection{Hard-Coded Venezuelan Locale}

\begin{badbox}[Venezuelan locale hard-coded throughout]
\texttt{America/Caracas} is hard-coded in the campaign form, the
storefront \texttt{link2Campaign}, and the \texttt{Cron} constructor
(v4). The subdomain regex is hard-coded to \texttt{necoyoad.com}
(v5). The tracking pixel is hard-coded to
\texttt{necoyoad.com/.../firma\_necoyoad.jpg} (v4). These make the
platform difficult to deploy for non-Venezuelan markets.
\end{badbox}

\subsection{\texttt{var\_dump()} in Production Code}

\begin{badbox}[var-dump left in the links widget]
The \texttt{ControllerModuleLinks::getLinks()} method (v7, line 56)
contains a \texttt{var\_dump(\$descriptions)} call that fires when
\texttt{submenu\_type === 'html\_content'}. Every menu link with an
HTML-content submenu dumps the descriptions array to the browser.
\end{badbox}

\subsection{No File Versioning in CSS/TPL Editor}

\begin{badbox}[CSS/TPL editor has no backup or versioning]
The \texttt{style/editor.php} controller (v6) lets the admin edit CSS
and TPL files directly with \texttt{file\_put\_contents()}. There is
no backup, no versioning, no diff. A typo in a CSS file can break the
storefront with no recovery path short of a server backup.
\end{badbox}

\subsection{Double Link Rewriting in Campaigns}

\begin{badbox}[Per-recipient link rewriting happens twice]
The campaign commit phase (v4) rewrites links and inserts
\texttt{campaign\_link} rows. The cron sender rewrites them
\emph{again} and inserts fresh rows with new nonces. The commit-phase
rows are dead data if the per-recipient cache hits. The
\texttt{campaign\_link} lookup is not scoped by
\texttt{campaign\_id}, so a nonce from one campaign can resolve to a
redirect from another.
\end{badbox}

% ============================================================
\part{What to Keep}\label{part:keep}
% ============================================================

\section{Patterns Worth Preserving in a Remastering}\label{sec:keep}

\subsection{The Declarative Admin CRUD Pattern}

\begin{reusebox}[AdminController declarative pattern $\to$ Laravel Form Requests + Resource Controllers]
The AdminController's declarative approach (\texttt{\$form\_vars},
\texttt{\$filters}, \texttt{\$public\_methods}) is worth preserving.
In Laravel, this maps to:
\begin{itemize}[leftmargin=1.6em,itemsep=2pt]
  \item \texttt{\$form\_vars} $\to$ Form Request validation rules
        (with type inference).
  \item \texttt{\$filters} $\to$ Query scopes on the Eloquent model.
  \item \texttt{\$public\_methods} $\to$ Middleware-based route
        permission checks.
  \item \texttt{\$model\_route} $\to$ Route model binding.
  \item \texttt{\$object\_type} $\to$ Eloquent morph map.
\end{itemize}
A base \texttt{AdminController} class could still provide
\texttt{index()}, \texttt{store()}, \texttt{update()},
\texttt{destroy()}, \texttt{grid()} --- but leveraging Laravel's
Form Requests, Resource Controllers, and Eloquent instead of raw
SQL.
\end{reusebox}

\subsection{The Polymorphic Object Spine}

\begin{reusebox}[Polymorphic spine $\to$ Laravel Polymorphic Relations]
The five polymorphic tables map directly to Laravel's polymorphic
relations:
\begin{itemize}[leftmargin=1.6em,itemsep=2pt]
  \item \texttt{description} $\to$ \texttt{MorphMany} on each model.
  \item \texttt{property} $\to$ \texttt{MorphMany} with a
        \texttt{group} + \texttt{key} composite.
  \item \texttt{object\_to\_store} $\to$ \texttt{MorphToMany}
        (Store model).
  \item \texttt{object\_to\_category} $\to$ \texttt{MorphToMany}
        (Category model).
  \item \texttt{url\_alias} $\to$ \texttt{MorphMany} with a
        \texttt{language\_id} scope.
\end{itemize}
The morph map would register each content type explicitly
(\texttt{'product' => Product::class}, etc.), which is cleaner than
the free-form string discriminator but loses some flexibility.
\end{reusebox}

\subsection{The Dual Hooks/Events System}

\begin{reusebox}[Hooks + Events $\to$ Laravel Events + Filters]
Laravel's event system maps to Necoyoad's Events (observer pattern).
For the Hooks (filter/action with short-circuit), Laravel doesn't
have a direct equivalent, but the pipeline pattern (middleware) or
a custom \texttt{Filter} facade could replicate it. Alternatively,
Laravel's \texttt{Event::dispatch()} with return values (Laravel 9+)
can serve as both observer and filter.
\end{reusebox}

\subsection{The Widget Composition System}

\begin{reusebox}[Widget composition $\to$ Laravel Blade Components + View Composers]
The widget composition system maps to Laravel's Blade components:
\begin{itemize}[leftmargin=1.6em,itemsep=2pt]
  \item Each widget becomes a Blade component
        (\texttt{<x-widgets.banner>} with settings prop).
  \item The \texttt{\{\%widget\_name\%\}} placeholder becomes a
        Blade slot or a dynamic component
        (\texttt{@stack('widgets')} or a dynamic component).
  \item The \texttt{loadWidgets()} method becomes a View Composer
        that populates the view with widget data.
  \item The per-entity widget override becomes a scoped View
        Composer that checks the session's \texttt{object\_type}.
  \item The visual editor's \texttt{nt-editable} attributes can be
        preserved with a Blade directive
        (\texttt{@editable('widget', \$name)}).
\end{itemize}
\end{reusebox}

\subsection{The Per-Entity Template Override}

\begin{reusebox}[Per-entity template override $\to$ Laravel dynamic view resolution]
Laravel's \texttt{View::exists()} + \texttt{view()} can replicate
the 3-level template resolution:
\begin{lstlisting}[language=PHP,basicstyle=\ttfamily\small,numbers=none,frame=none]
$template = $entity->properties['style']['view'] ?? null;
$default = config("defaults.{$entity->type}", "content/{$entity->type}");
$template = $template ?: $default;
if (view()->exists("themes/{$theme}/{$template}")) {
    return view("themes/{$theme}/{$template}", $data);
} else {
    return view("themes/choroni/{$template}", $data);
}
\end{lstlisting}
\end{reusebox}

\subsection{The Multi-Store/Multi-Language Architecture}

\begin{reusebox}[Multi-store $\to$ Laravel multi-tenancy + localization]
Laravel has mature multi-tenancy packages (e.g.\
\texttt{tenancy/tenancy}, \texttt{stancl/tenancy}) that provide
database-level or row-level multi-tenancy. The row-level approach
(\texttt{store\_id} on every table) maps to a global scope that
filters by the current tenant. Laravel's built-in localization
(\texttt{App::setLocale()}, translatable packages like
\texttt{spatie/laravel-translatable}) handles the language dimension.
\end{reusebox}

\subsection{The Campaign Marketing Pipeline}

\begin{reusebox}[Campaign pipeline $\to$ Laravel Queues + Mail + Horizon]
The campaign pipeline maps cleanly to Laravel's queue infrastructure:
\begin{itemize}[leftmargin=1.6em,itemsep=2pt]
  \item \texttt{task\_queue} $\to$ Laravel Queue (Redis/Database driver).
  \item \texttt{CronSend::sendCampaign()} $\to$ a queued Job
        (\texttt{SendCampaignEmail}).
  \item \texttt{task\_exec} locking $\to$ Laravel's built-in job
        locking (\texttt{ShouldBeUnique}).
  \item PHPMailer 5.0 $\to$ Laravel Mail (Symfony Mailer under the
        hood, with OAuth2 support).
  \item 50-email throttle $\to$ Laravel's \texttt{RateLimited} middleware.
  \item Tracking pixel + link redirect $\to$ controller routes
        with \texttt{DB::insert()} on each hit.
  \item Bounce processing $\to$ Laravel's
        \texttt{Mail::failure()} callbacks + a bounce-handling job.
\end{itemize}
\end{reusebox}

\subsection{The Banner jQuery Plugin Discriminator}

\begin{reusebox}[jquery-plugin discriminator $\to$ Laravel dynamic Blade components]
The \texttt{jquery\_plugin} discriminator maps to Laravel's dynamic
Blade components:
\texttt{<x-dynamic-component>} with a component prop. Each slider becomes a Blade component
with a component prop. Each slider becomes a Blade component
(\texttt{resources/views/components/sliders/nivo-slider.blade.php}).
The \texttt{ntPlugins} registry can be replaced by Alpine.js or
Livewire for client-side initialisation.
\end{reusebox}

% ============================================================
\part{Feasibility of a Laravel Rebuild}\label{part:laravel}
% ============================================================

\section{Pattern-by-Pattern Mapping}\label{sec:mapping}

\begin{table}[htbp]
\centering
\caption{Necoyoad $\to$ Laravel pattern mapping.}\label{tab:laravel-mapping}
\footnotesize\sloppy
\begin{tabularx}{\columnwidth}{>{\raggedright\arraybackslash}p{4.0cm}>{\raggedright\arraybackslash}p{4.5cm}X}
\toprule
\textbf{Necoyoad pattern} & \textbf{Laravel equivalent} & \textbf{Effort} \\
\midrule
Front Controller + Action & Laravel Routing + Controllers & Low --- Laravel's router is more capable. \\
Registry (service locator) & Laravel Service Container & Low --- automatic DI. \\
AdminController declarative CRUD & Form Requests + Resource Controllers & Medium --- needs a base controller. \\
Polymorphic object spine & Eloquent Morph Relations & Medium --- morph map per type. \\
Hooks (WordPress-style) & Custom Filter facade or Pipeline & Medium --- no built-in equivalent. \\
Events (Node-style) & Laravel Events + Listeners & Low --- direct mapping. \\
Widget composition & Blade Components + View Composers & High --- core feature, needs careful design. \\
\texttt{\{\%widget\_name\%\}} substitution & Blade slots / dynamic components & Medium. \\
Per-entity template override & \texttt{view()->exists()} + dynamic resolution & Low. \\
Multi-store (data-level) & Multi-tenancy package + global scope & Medium --- tenant resolution middleware. \\
Multi-language (6-level priority) & Laravel Localization + middleware & Low. \\
Campaign pipeline & Laravel Queues + Mail + Horizon & Medium --- job design + tracking. \\
Banner jquery\_plugin & Dynamic Blade components & Low. \\
ntPlugins JS registry & Alpine.js / Livewire / JS modules & Medium --- front-end rewrite. \\
SEO URL rewriter & Route Model Binding + slug column & Low. \\
MyISAM tables & MySQL InnoDB + migrations & Medium --- schema migration. \\
Double-MD5 passwords & bcrypt / argon2id & Low --- migration hash check. \\
Cron CLI & Laravel Scheduler + Queue Worker & Low. \\
Visual theme editor & Livewire / Inertia + contenteditable & High --- significant front-end work. \\
deps.php asset manifest & Laravel Vite + manifest & Low. \\
\bottomrule
\end{tabularx}
\end{table}

\section{Effort Estimate}\label{sec:effort}

\begin{table}[htbp]
\centering
\caption{Estimated effort by subsystem (senior developer-months).}\label{tab:effort}
\small
\begin{tabularx}{\columnwidth}{lrX}
\toprule
\textbf{Subsystem} & \textbf{Effort} & \textbf{Notes} \\
\midrule
Core framework (routing, DI, auth) & 1 mo & Mostly leveraging Laravel builtins. \\
Database migration (InnoDB + schema) & 2 mo & 87 tables, data migration, FK constraints. \\
Admin CRUD framework & 2 mo & Base controller + Form Requests + 70+ entities. \\
Widget composition system & 3 mo & The hardest part --- Blade components, view composers, visual editor. \\
Storefront (catalog, checkout, account) & 3 mo & Product, category, cart, order, customer. \\
CMS (posts, pages, categories) & 1 mo & Simpler than e-commerce. \\
Marketing pipeline & 2 mo & Queue jobs, mail, tracking, bounce. \\
Banner subsystem & 1 mo & Dynamic components, 33 slider templates. \\
Multi-store / multi-language & 1 mo & Tenancy package + localization. \\
Visual theme editor & 3 mo & Significant front-end (Livewire/Inertia + drag-and-drop). \\
API (REST, with proper auth) & 1 mo & Laravel Sanctum + API Resources. \\
Testing + deployment & 2 mo & PHPUnit, CI/CD, Docker. \\
\midrule
\textbf{Total} & \textbf{22 mo} & $\sim$2 developer-years (senior). \\
\bottomrule
\end{tabularx}
\end{table}

\section{Recommended Migration Strategy}\label{sec:strategy}

\subsection{Option A: Full Rebuild with Laravel (Recommended)}

\begin{enumerate}[leftmargin=1.6em,itemsep=2pt]
  \item Start with a fresh Laravel 11 project.
  \item Migrate the database to InnoDB with proper foreign keys.
  \item Implement the polymorphic object spine as Eloquent traits.
  \item Build the AdminController base + Form Request pattern.
  \item Port the 70+ admin entities (mostly declarative, should be fast).
  \item Build the widget composition system as Blade components + View
        Composers (the hardest part).
  \item Port the storefront (catalog, checkout, account).
  \item Port the CMS (posts, pages).
  \item Port the marketing pipeline as Queue jobs.
  \item Port the banner subsystem as dynamic Blade components.
  \item Build the visual theme editor with Livewire/Inertia.
  \item Build the API with Sanctum + API Resources.
  \item Deploy with Docker + CI/CD.
\end{enumerate}

\textbf{Pros}: clean codebase, modern framework, proper testing,
security best practices, active community.
\textbf{Cons}: 22 developer-months of effort, no incremental revenue
during the rebuild.

\subsection{Option B: Incremental Modernisation (Strangler Fig)}

\begin{enumerate}[leftmargin=1.6em,itemsep=2pt]
  \item Keep the existing codebase running.
  \item Add Laravel alongside (via \texttt{require} in
        \texttt{composer.json}).
  \item Migrate one subsystem at a time (start with the API, then
        admin, then storefront).
  \item Use Laravel's database connection to access the existing
        MyISAM tables (read-only initially, then read-write as each
        subsystem is migrated).
  \item Migrate the database to InnoDB table-by-table.
  \item Eventually decommission the old codebase.
\end{enumerate}

\textbf{Pros}: incremental, no big-bang risk, revenue continues.
\textbf{Cons}: longer total timeline ($\sim$30 months), complexity
of running two codebases.

\subsection{Option C: Remaster in Place (Modernise PHP, Keep Architecture)}

\begin{enumerate}[leftmargin=1.6em,itemsep=2pt]
  \item Migrate MyISAM $\to$ InnoDB.
  \item Replace double-MD5 with bcrypt.
  \item Fix the \texttt{object\_to\_store} inconsistencies.
  \   Fix the broken API PUT branches.
  \item Add unsubscribe mechanism.
  \item Replace PHPMailer 5.0 with PHPMailer 6.x.
  \item Fix the \texttt{var\_dump} and other bugs.
  \item Add PHPUnit tests.
\end{enumerate}

\textbf{Pros}: lowest effort ($\sim$6 developer-months), keeps the
working codebase.
\textbf{Cons}: doesn't address the fundamental architecture issues
(MyISAM, no DI container, raw SQL, no testing culture).

\section{Recommendation}\label{sec:recommendation}

\textbf{Option A (Full Rebuild with Laravel) is recommended} if the
goal is a modern, maintainable platform that can be extended for the
next 10 years. The 22-month effort is significant but the result is
a clean codebase with proper testing, security, and community support.

\textbf{Option B (Strangler Fig) is recommended} if there is a
running business that cannot afford a big-bang rebuild. The
incremental approach reduces risk and allows revenue to continue,
at the cost of a longer timeline and the complexity of running two
codebases.

\textbf{Option C (Remaster in Place) is recommended} as a short-term
stabilisation step before Option A or B. The 6-month effort fixes
the most critical security and data-integrity issues without
disrupting the business.

\section{What Would \emph{Not} Survive a Rebuild}\label{sec:discard}

\begin{itemize}[leftmargin=1.6em,itemsep=2pt]
  \item The raw SQL query builder (\texttt{Model::buildSQLQuery()}) ---
        replaced by Eloquent.
  \item The \texttt{LIKE '\%key=value\%'} widget query pattern ---
        replaced by proper JSON columns or a junction table.
  \item The \texttt{\{\%widget\_name\%\}} string substitution ---
        replaced by Blade slots.
  \item The MyISAM tables --- replaced by InnoDB.
  \item The double-MD5 password hash --- replaced by bcrypt.
  \item The \texttt{validateTokens()} stub --- replaced by Sanctum.
  \item The PHPMailer 5.0 wrapper --- replaced by Laravel Mail.
  \item The \texttt{.htaccess}-based routing --- replaced by Laravel's
        router.
  \item The \texttt{map.php} 843-line switch statement --- replaced by
        Laravel's autoloading + service provider pattern.
  \item The hand-rolled session handler --- replaced by Laravel's
        session driver.
  \item The file-based cache --- replaced by Redis.
\end{itemize}

% ============================================================
\section{Conclusion}\label{sec:conclusion}
% ============================================================

Necoyoad is a characterful, ambitious platform that was clearly built
with care and a strong vision. The AdminController declarative CRUD
framework, the polymorphic object spine, the dual Hooks/Events system,
and the widget composition system are genuinely innovative --- they
anticipate patterns that modern frameworks are only now adopting
(Laravel's dynamic components, WordPress's block editor, Drupal's
Layout Builder).

The technical debt (MyISAM, double-MD5, raw SQL, no testing) is
significant but not unusual for a platform of this vintage (2009--2023).
The bugs (broken API PUTs, vestigial bounce, var\_dump in production,
missing store\_id) are fixable but indicative of a codebase that grew
faster than its testing infrastructure.

A Laravel rebuild is feasible and recommended. The architecture maps
well: the polymorphic spine maps to Eloquent morph relations, the
widget system maps to Blade components, the AdminController maps to
Form Requests, the marketing pipeline maps to Queue jobs. The 22-month
effort estimate is for a senior developer; a team of 2--3 could
complete it in 8--12 months.

The result would be a modern, testable, secure e-commerce + CMS
platform that preserves the best of Necoyoad's architecture while
leaving behind its technical debt.

\end{document}
